\clearpage
\fancyhead[L]{\small BESTIMMUNG DER ABSOLUTEN DEKLINATION.}
\SourceLeaf{442}\pdfbookmark[0]{Über die Anwendung des Magnetometers zur Bestimmung der absoluten Deklination.}{absolute-declination-title}
\begin{center}
{\Large\bfseries ÜBER DIE ANWENDUNG DES MAGNETOMETERS\par\medskip ZUR BESTIMMUNG DER ABSOLUTEN DEKLINATION.}
\par\bigskip
Resultate aus den Beobachtungen des magnetischen Vereins. 1841. I.
\end{center}
Ich beabsichtige nicht, den in der Überschrift genannten Gegenstand hier noch einmal vollständig zu behandeln. Darüber ist bereits in Band 2 der \textit{Resultate} [1837. VII.] ein sehr ausführlicher Aufsatz erschienen. Ich werde mich hier auf eine einzige Hauptaufgabe beschränken, für die die dort auf S. 121—124 gegebene Herleitung unzureichend erscheint. Es geht um die Bestimmung des Azimuts der Vertikalebene, in der die optische Achse des Beobachtungsfernrohrs liegt.

Die betreffende Vertikalebene ist durch die Marke und einen festen Punkt der Skala festgelegt. Dieser Punkt wird durch den Lotfaden bestimmt, der über der Mitte des Objektivs des Beobachtungsfernrohrs herabhängt. Vom Standort des Beobachtungsfernrohrs aus muss ein entfernter Gegenstand sichtbar sein, dessen Azimut bereits anderweitig bekannt ist. Zunächst ist also der auf die Horizontalebene projizierte Winkel zwischen diesem Gegenstand und der Marke zu bestimmen. Ich nehme an, dass dafür ein Theodolit der bekannten, von \textsc{Reichenbach} eingeführten Bauart verwendet wird. Damit setze ich nicht voraus, dass derselbe Theodolit auch für die magnetischen Beobachtungen verwendet wird. Dafür kann ohne Weiteres ein besonderes Ablesefernrohr dienen.

Solche Theodolite werden gewöhnlich zu Winkel\-messungen zwischen Gegenständen in so großer Entfernung verwendet, dass eine geringe
\expandafter\def\csname HistoricalNoteText443\endcsname{Bei den weiter unten anzuführenden Beobachtungen etwa 20 Millimeter.}%
\SourceLeaf{443}Abweichung bei der Ausführung von mehreren Bedingungen, die dem Entwurf des Instruments zugrunde liegen, keinen merklichen Fehler hervorrufen kann. Tatsächlich lässt sich bei keiner mechanischen Arbeit absolute Vollkommenheit erreichen. Wenn die Gegenstände aber vergleichsweise sehr nahe sind (oder, wie im vorliegenden Fall, einer von ihnen), muss man solche Abweichungen allerdings genauer beachten. Insbesondere sind hier die folgenden Umstände zu berücksichtigen.

I. Die vertikale Drehachse, die horizontale Drehachse und die optische Achse des Fernrohrs sollten sich in einem einzigen Punkt schneiden. Soweit diese Bedingung nicht vollständig erfüllt ist, tritt eine dreifache Abweichung auf. Es seien $A,B$ die beiden Punkte auf der vertikalen beziehungsweise der horizontalen Drehachse, an denen die Achsen einander am nächsten sind. Ebenso seien $C,D$ die entsprechenden Punkte auf der horizontalen Drehachse und der optischen Achse. Man bezeichne die Abstände $AB,BC,CD$ mit $\alpha,\beta,\gamma$ und lege die Vorzeichen nach Belieben fest. Beispielsweise kann man $\alpha$ als positiv ansetzen, wenn $A$ auf derselben Seite der horizontalen Drehachse liegt wie das Okular des Fernrohrs, $\beta$ als positiv, wenn für den Beobachter am Okular $C$ rechts von $B$ liegt, und $\gamma$ als positiv, wenn $D$ oberhalb von $C$ liegt.

II. Die optische Achse des Fernrohrs sollte senkrecht zur horizontalen Drehachse stehen. Diese Bedingung lässt sich zwar mit aller erforderlichen Genauigkeit erfüllen. Um aber nach der Beobachtung eines entfernten Gegenstands einen nahen deutlich sehen zu können, muss man das Okularrohr weiter\HistoricalNote{443} herausziehen und dadurch dem Fadenkreuz eine andere Stellung zum Objektiv geben. Man darf deshalb nicht voraussetzen, dass beiden Stellungen des Okularrohrs dieselbe optische Achse entspricht, und muss damit rechnen, dass die für eine Stellung vorgenommene Berichtigung in der anderen wieder verloren geht. Der größeren Allgemeinheit halber kann man annehmen, dass die Berichtigung für keine der beiden Stellungen genau vorgenommen wurde, und den Kollimationsfehler für die erste Stellung mit $c$, für die zweite mit $c'$ bezeichnen. Man kann die Fehler als positiv ansetzen, wenn die optische Achse mit dem rechts vom Beobachter liegenden Arm der horizontalen Achse einen spitzen Winkel bildet.

\SourceLeaf{444}Gehören die Größen $\beta,\gamma$ zur ersten Okularstellung, so treten bei der zweiten offenbar etwas veränderte Werte an ihre Stelle. Wir können sie mit $\beta',\gamma'$ bezeichnen.

Zwar lässt sich der Einfluss aller dieser Abweichungen auf die Messung sowohl des horizontalen Winkels zwischen den beiden Gegenständen als auch ihrer Höhenwinkel (wenn der Theodolit zugleich einen Höhenkreis hat) leicht in exakten Gleichungen darstellen, aus denen sich die Ergebnisse mithilfe einer Gleichung vierten Grades ableiten ließen. Da aber die sieben Größen $\alpha,\beta,\gamma,\beta',\gamma',c,c'$ alle nur sehr klein sein können, kann man unbedenklich alle Größen. die in Bezug auf diese Größen zweiter oder höherer Ordnung sind, vernachlässigen und das Ergebnis ihres Einflusses in eine sehr einfache Form bringen. Hier können wir aber selbst auf diese Darstellung verzichten. Man sieht nämlich leicht, dass alle sieben Abweichungen lediglich das entgegengesetzte Vorzeichen annehmen, ohne ihren Betrag zu ändern, wenn man das Fernrohr auf die richtige Weise umlegt. Dasselbe gilt demnach auch für die Fehler der Messungen, die man bei den beiden verschiedenen Arten der Lagerung vornimmt. Der Mittelwert dieser beiden Messungen ist folglich von selbst vom Einfluss dieser Fehler frei, ohne dass man ihre einzelnen Anteile zu kennen braucht. So erhält man den wahren Wert des Winkels zwischen den beiden Vertikalebenen, die sich in der Vertikalachse des Instruments schneiden und in denen die beiden Gegenstände liegen. Dasselbe gilt für die Höhenwinkel, die sich dann auf den Punkt $A$ beziehen, für unseren gegenwärtigen Zweck aber nicht nötig sind.

Das Umlegen muss so geschehen, dass die Zapfen wieder in dieselben Lagerpfannen kommen, während die obere Seite des Fernrohrs zur unteren wird und das Objektiv an die Stelle des Okulars tritt. Dieses Umlegen entspricht also einer halben Umdrehung um die horizontale Achse, für deren Ausführung die Stützen lediglich nicht hoch genug sind. Wollte man das Fernrohr stattdessen so umlegen, dass die Zapfen in die jeweils anderen Lagerpfannen gelegt werden, während das Objektivende auf derselben Seite bleibt (was geometrisch einer halben Umdrehung um die Achse des Fernrohrs entspricht), dann würden nicht alle sieben Größen $\alpha,\beta,\gamma,\beta',\gamma',c,c'$ einfach das entgegengesetzte Vorzeichen annehmen. Nur $\gamma,\gamma',c,c'$ würden einfach das entgegengesetzte Vorzeichen annehmen. Es ist nämlich nicht sicher, dass die Stützen \textit{genau}
\expandafter\def\csname HistoricalNoteText445\endcsname{Auf der ersten Tafel des ersten Bandes der \textit{Resultate} [1836. I.] ist dieser Turm angedeutet, ungefähr so, wie er bei geschlossenem Fenster vom Standort des Theodoliten aus erscheint. Von dem Augenpunkt aus, der der perspektivischen Zeichnung eigentlich zugrunde liegt, wird der Turm durch die Wand links vom Fenster verdeckt.}%
\SourceLeaf{445}gleich weit von der Vertikalachse entfernt sind, und daher. könnte der Punkt $B$ nach einem solchen Umlegen ein anderer sein als vorher, sodass auch $\beta$ und $\beta'$ andere Werte annehmen könnten. Dass $\alpha$ dabei das entgegengesetzte Vorzeichen nicht annimmt und genau den vorigen Wert behält, ist hier allerdings unwesentlich, weil $\alpha$ im linearen Ausdruck für den Fehler der horizontalen Winkelmessung überhaupt nicht vorkommt.

Wie sich eine solche Winkelmessung für den beabsichtigten Zweck nutzen lässt, zeigt sich am einfachsten an einem Beispiel. Dafür wähle ich die letzte Anwendung des Verfahrens, die am 11. März dieses Jahres ausgeführt wurde.

Im hiesigen magnetischen Observatorium dient ein Stadtkirchturm dazu, die Beobachtungen an den wahren astronomischen Meridian anzuschließen. Seine Knopfstange ist am Standort des Beobachtungsfernrohrs durch das geöffnete nordliche Fenster frei sichtbar\HistoricalNote{445}, und zwar von der Mitte der Säule aus, die seit Juli 1837 an der Stelle des früher verwendeten hölzernen Stativs steht, im Azimut $173^0 35'25''5$. Dieses Azimut wurde bestimmt, indem man an einer anderen Stelle des Saals einen Theodoliten aufstellte, dessen Vertikalachse genau in der Flucht der Säulenmitte und des Kirchturms lag, und die Winkel zwischen diesem und zwei weiteren von dort sichtbaren Kirchtürmen maß. Die Lage dieser Türme zum Nullpunkt in der Sternwarte war durch frühere, an die Gradmessung angeschlossene Messungen genau bekannt, und das betreffende Azimut ließ sich daher aus jenen Winkelmessungen leicht berechnen.

Nun wurde ein achtzölliger \textsc{Ertel}scher Repetitionstheodolit so auf der Säule aufgestellt, dass seine Vertikalachse möglichst genau durch die Mitte der Säule ging. Der horizontale Winkel zwischen der Marke und der Knopfstange des Turms wurde bei den beiden oben beschriebenen Arten der Lagerung des Fernrohrs durch jeweils 25 Repetitionen gemessen. In der ersten Lage ergab sich der Winkel
\[
=11^0 40'54''50
\]
\SourceLeaf{446}in der zweiten
\[
=11^0 41'36''18
\]

Der wahre Wert des Winkels mit dem Scheitel in der Vertikalachse des Theodoliten ist folglich
\[
=11^0 41'15''34
\]
und damit das Azimut der Vertikalebene durch diese Vertikalachse und die Marke
\[
=161^0 54'10''16
\]

Mit diesem Vorgang war ein weiterer verbunden. Er sollte ergeben, in welchem Punkt diese Vertikalebene die Skala schneidet.

Auf das Objektivende des Theodolitfernrohrs ist ein Ring gesteckt. Er hat auf seiner Vorderfläche zwei feine, einander diametral gegenüberliegende Einschnitte und auf seiner runden Außenfläche zwei entsprechende Häkchen. Je nach Lage des Fernrohrs wird ein feiner, mit einem Gewicht beschwerter Goldfaden in das eine oder andere Häkchen eingehängt. Man dreht den Ring so, dass der Durchmesser durch die Einschnitte senkrecht zur horizontalen Drehachse des Fernrohrs steht. Man erkennt, dass diese Bedingung erfüllt ist, wenn der im oberen Einschnitt liegende Lotfaden zugleich genau auf den unteren ausgerichtet ist. Dazu muss man das Fernrohr nahezu horizontal stellen und es nur so weit nach unten neigen, dass der Faden gerade noch frei vor dem Ring hängt. Die Übereinstimmung prüft man mit einer Lupe. Der Lotfaden hängt in sehr geringer Entfernung vor der Skala. Nun sind die entsprechenden Punkte der Skala in den beiden verschiedenen Lagen des Fernrohrs zu notieren, wobei es jedes Mal auf die Marke oder genauer in deren Vertikalebene gerichtet wird. Genau genommen sind damit die Punkte der Skala gemeint, die in der Vertikalebene durch die Marke und den Lotfaden liegen. Das kann man unmittelbar im Spiegel des Magnetometers erkennen, wenn der Theodolit selbst als Ablesefernrohr dienen soll, die Skala sich also in einer dafür geeigneten Höhe befindet. Es ist wohl überflüssig darauf hinzuweisen, dass es in diesem Fall notwendig werden kann, den Magnetstab
\SourceLeaf{447}des Magnetometers mithilfe eines aus der Ferne wirkenden Ablenkungsstabs erst in eine Stellung zu bringen, in der der betreffende Skalenpunkt nahe am Fadenkreuz des Theodolitfernrohrs erscheint. Im hiesigen magnetischen Observatorium werden die magnetischen Beobachtungen jetzt mit einem besonderen Ablesefernrohr vorgenommen, das sich weniger hoch über der Säule befindet als das Theodolitfernrohr. Mit diesem ist das Bild der Skala im Spiegel des Magnetometers nicht sichtbar, da die Skala in einer der Lage des Ablesefernrohrs entsprechenden Höhe angebracht ist. Zur Bestimmung des Skalenpunkts, der dem vom Theodolitfernrohr herabhängenden Lotfaden entspricht, habe ich daher das Ablesefernrohr selbst verwendet. Es wurde zu diesem Zweck nahe der Marke in der betreffenden Vertikalebene aufgestellt. Dass man es mit dieser letzten Bedingung nicht allzu genau nehmen muss, versteht sich von selbst, da der Lotfaden, wie bereits bemerkt, in geringer Entfernung vor der Skala hängt. Auf diese Weise fand man, dass der Lotfaden
\begin{quote}
bei der ersten Lage des Theodolitfernrohrs dem Skalenpunkt $850.8$ entsprach und

 bei der zweiten Lage dem Punkt $849.4$,
\end{quote}
woraus man schließen darf, dass die Vertikalebene durch die Marke und die Vertikalachse des Theodoliten, deren Azimut oben bestimmt wurde, die Skala im Punkt $850.1$ schneidet.

Die Bestimmung des Azimuts der Vertikalebene, in der die optische Achse des Beobachtungsfernrohrs liegt, bietet nun keine weitere Schwierigkeit. Entspricht der vor der Mitte seines Objektivs herabhängende Lotfaden dem Skalenpunkt $850{,}1+n$, so genügt es (da vorausgesetzt wird, dass die Skala senkrecht zu jener Ebene steht), das Produkt $n.206265''$ durch die horizontale Entfernung der Skala von der Marke, ausgedrückt in Skalenteilen, zu teilen und den Quotienten mit seinem Vorzeichen zu $161^0 54'10''16$ zu addieren. Gegenwärtig ist diese Entfernung $=9638.7$. Diente also der Theodolit selbst zum Beobachten, so wäre dieses Azimut in der ersten Lage des Fernrohrs
\[
=161^0 54'25''1
\]
in der zweiten Lage dagegen
\[
=161^0 53'55''2
\]
\SourceLeaf{448}Zum Beobachten dient aber, wie bereits bemerkt, ein besonderes Ablesefernrohr. Dieses wurde nach Abschluss der obigen Vorgänge so aufgestellt, dass der Lotfaden dem Punkt $850.0$ entsprach, wenn die optische Achse auf die Vertikale durch die Marke gerichtet war. Das gesuchte Azimut ist daher
\[
=161^0 54'8''0
\]

Zu dem hier behandelten Verfahren seien noch ein paar Bemerkungen hinzugefügt.

I. Hat die horizontale Achse in ihren Lagern etwas Spiel in Längsrichtung, muss man dafür sorgen, dass sie bei den einzelnen Winkelmessungen immer dieselbe Lage zu den Stützen hat. Dazu kann man jedes Mal durch leichten Druck gegen das Ende eines bestimmten Zapfens das Spiel auf einer Seite beseitigen. Ohne diese Vorsicht könnte man nicht darauf rechnen, dass die oben mit $\beta$ bezeichnete Größe in der ersten Lage des Fernrohrs bei allen Repetitionen immer denselben Wert und in der zweiten immer genau den entgegengesetzten Wert behält.

II. Für die hier beschriebenen Vorgänge braucht die optische Achse des Theodolitfernrohrs in keiner der beiden Okularstellungen genau justiert zu sein, also senkrecht zur horizontalen Drehachse zu stehen. Dient der Theodolit aber zugleich als Ablesefernrohr, muss diese Justierung vor einem solchen Gebrauch tatsächlich vorgenommen werden, und zwar für die Stellung des Okularrohrs, in der man beobachtet oder in der die Marke und das Spiegelbild der Skala deutlich erscheinen. Bekanntlich prüft man durch Umlegen, ob die optische Achse senkrecht zur horizontalen Drehachse steht. Man benutzt dazu gerade das Umlegen, das bei den obigen Winkelmessungen \textit{nicht} angewandt werden durfte [S. 439 d. B.], indem man die Zapfen in die jeweils gegenüberliegenden Lager legt, ohne die Richtung, in die das Fernrohr zeigt, umzukehren. Gewöhnlich bezieht sich eine solche Prüfung auf die Stellung des Okularrohrs, in der man sehr weit entfernte Gegenstände deutlich sieht. und in diesem Fall ist tatsächlich nichts weiter nötig, als dass ein solcher Gegenstand vor und nach dem Umlegen auf dem Fadenkreuz erscheint. Wenn aber im vorliegenden Fall der vor der Mitte des Objektivs herabhängende Lotfaden nach dem Umlegen eine andere Lage hat als zuvor, muss man neben dem ersten Zielpunkt einen zweiten verwenden, der um denselben Betrag versetzt ist. Offenbar muss auch ein an%
\SourceLeaf{449}%
stelle des Theodoliten verwendetes besonderes Ablesefernrohr derselben Justierung unterzogen werden und dafür geeignet aufgestellt sein. Selbstverständlich muss auch die horizontale Drehachse richtig nivelliert sein. Die beiden bei den hiesigen Magnetometern verwendeten Ablesefernrohre sind fast genauso aufgestellt wie Theodolite, nur ohne geteilte Kreise, und haben eine erheblich stärkere optische Leistung als Theodolitfernrohre gewöhnlich haben.

Es sei noch bemerkt, dass der \textit{Einfluss} eines Kollimationsfehlers auf das Azimut der optischen Achse nur ein sehr kleiner Bruchteil des Kollimationsfehlers selbst ist. Dieser Bruchteil wird durch die Differenz der Sekanten der beiden Neigungen bestimmt, wenn das Fernrohr einmal auf die Marke und dann auf den Spiegel gerichtet ist. Beim hiesigen Unifilarmagnetometer betragen diese Neigungen $1^0 55'$ und $5^0 16'$. Der Unterschied der Azimute der optischen Achse bei der Ausrichtung auf Marke und Spiegel beträgt folglich nur $\frac{1}{272}$ des Kollimationsfehlers selbst. Unter ähnlichen Umständen wird es daher gewöhnlich genügen, die Kollimation an einem entfernten Gegenstand zu justieren. Denn sofern das Fadenkreuz durch diese Justierung nicht weit aus der Mitte des Okularrohrs geraten ist, kann ein weiteres Herausziehen des Okularrohrs kaum einen Kollimationsfehler von mehr als einem kleinen Bruchteil einer Bogenminute erzeugen, sodass dessen Einfluss ganz unmerklich bleibt.

III. Der Lotfaden bleibt ständig am Beobachtungsfernrohr hängen, damit eine unbeabsichtigte Verschiebung der Skala sofort erkennbar wird. Tritt eine solche Verschiebung ein, kann man die Skala entweder wieder in ihre frühere Stellung bringen oder in der Rechnung von dem Punkt der Skala aus zählen, der dem Lotfaden nach der Verschiebung entspricht, und zwar ebenso wie zuvor vom früheren Punkt aus. Bei der gegenwärtigen Art, die Skala im magnetischen Observatorium an der Säule zu befestigen, kommen unbeabsichtigte Verschiebungen allerdings überhaupt nicht mehr vor.

